\section{Problem\'atica de la volatilidad de precios agr\'icolas}

\subsection{Comportamiento variable de los precios agr\'icolas}

Los precios de productos agr\'icolas de exportaci\'on, incluida la palta Hass, exhiben comportamientos variables en el tiempo: tendencias no lineales, ciclos marcados, cambios de r\'egimen y eventos extremos \citep{Velasquez2025, Zhang2025}. En t\'erminos econom\'etricos, muchas series de precios agr\'icolas son no estacionarias, es decir, presentan media, varianza o estructura de autocorrelaci\'on que cambia con el tiempo, lo que dificulta extrapolar tendencias simples sin transformaciones adecuadas \citep{Xia2024}.

\citet{VitteryFlores2023} documentaron una ca\'ida del precio promedio FOB peruano desde 2.38 USD/kg en 2016 hasta 1.53 USD/kg en 2022, seguida de recuperaciones parciales en a\~nos posteriores. Esta trayectoria no responde a un \'unico factor, sino a la interacci\'on entre mayor oferta mundial, competencia internacional, concentraci\'on temporal de despachos y variaciones en la demanda de mercados clave. La Figura~\ref{fig:valores_unitarios} muestra la evoluci\'on reciente de valores unitarios medios de exportaci\'on mundial de aguacate, evidenciando fluctuaciones mensuales relevantes para la planificaci\'on comercial.

\incluirfigura{figuras/figura2\_1.jpg}{0.92\linewidth}{Valores unitarios medios de exportaci\'on mundial de aguacate (ene 2023 -- jul 2025)}{fig:valores_unitarios}

La volatilidad no se limita al mercado peruano. \citet{Trincado2022} analiza precios en el mercado mayorista Lo Valledor (Chile), mientras \citet{FAO2026} reporta variaciones en precios al por mayor en Estados Unidos. La comparaci\'on entre mercados subraya que los productores enfrentan se\~nales de precio dispersas y, frecuentemente, desactualizadas respecto a las condiciones FOB que determinan su rentabilidad exportadora.

\subsection{Influencia de oferta y demanda, clima y estacionalidad}

La formaci\'on del precio FOB sintetiza fuerzas de oferta (hect\'areas sembradas, rendimiento, calibres disponibles), demanda (consumo en destino, inventarios, preferencias), estacionalidad (calendarios de cosecha por pa\'is) y perturbaciones clim\'aticas \citep{SierraSelva2020, Velasquez2025}. La Tabla~\ref{tab:estacionalidad_export} resume la estacionalidad de exportaciones mundiales, mostrando c\'omo M\'exico, Per\'u, Chile y Kenia concentran despachos en meses distintos pero parcialmente superpuestos.

\begin{table}[H]
\centering
\caption{Estacionalidad de las exportaciones mundiales de palta Hass}
\label{tab:estacionalidad_export}
\small
\begin{tabular}{lrrrr}
\toprule
\textbf{Mes} & \textbf{M\'exico (\%)} & \textbf{Per\'u (\%)} & \textbf{Chile (\%)} & \textbf{Kenia (\%)} \\
\midrule
Ene--Mar & 18 &  2 &  5 & 12 \\
Abr--Jun & 22 & 45 & 15 & 28 \\
Jul--Sep & 35 & 48 & 55 & 35 \\
Oct--Dic & 25 &  5 & 25 & 25 \\
\bottomrule
\end{tabular}
\parbox{0.94\textwidth}{\small \textit{Nota.} Elaboraci\'on propia con base en \citep{SierraSelva2020}, Cuadro No.~33. Los porcentajes representan la participaci\'on mensual aproximada en el total anual exportado.}
\end{table}


\begin{table}[H]
\centering
\caption{Periodos productivos de palta variedad Hass en Chile y Per\'u}
\label{tab:periodos_hass}
\small
\begin{tabular}{lll}
\toprule
\textbf{Pa\'is} & \textbf{Periodo principal de cosecha} & \textbf{Ventana exportadora clave} \\
\midrule
Chile & Septiembre--Abril & Octubre--Marzo \\
Per\'u & Marzo--Agosto & Abril--Julio \\
\bottomrule
\end{tabular}
\parbox{0.94\textwidth}{\small \textit{Nota.} Elaboraci\'on propia con base en \citep{Trincado2022}, Cuadro 2.3. La complementariedad estacional entre ambos pa\'ises condiciona la formaci\'on de precios internacionales.}
\end{table}


La Tabla~\ref{tab:periodos_hass} contrasta los periodos productivos de Chile y Per\'u, evidenciando una complementariedad parcial que no elimina competencia en transiciones de campa\~na. Eventos clim\'aticos --- sequ\'ias, heladas, exceso de lluvias --- modifican calendarios de cosecha y calidad de fruta, alterando la relaci\'on oferta--precio en semanas cr\'iticas \citep{Capunay2025}. Asimismo, la concentraci\'on de exportaciones en pocas semanas puede deprimir precios promedio aun cuando el volumen anual crezca.

Los costos log\'isticos, el tipo de cambio y los aranceles o barreras fitosanitarias a\~naden capas de incertidumbre. \citet{ComexPeru2025} se\~nala que la diversificaci\'on de mercados es estrat\'egica, pero no neutraliza la volatilidad cuando m\'ultiples or\'igenes compiten simult\'aneamente en los mismos destinos. Para el productor, la estacionalidad implica que decisiones de siembra tomadas a\~nos vista deben anticipar no solo el clima local, sino la configuraci\'on global de oferta en el momento de su primera cosecha comercial.

\subsection{Impacto econ\'omico en productores y exportadores}

La volatilidad de precios impacta directamente m\'argenes de ganancia, capacidad de reinversi\'on y estabilidad financiera de productores y exportadores \citep{Capunay2025, VargasCordova2025}. \citet{VargasCordova2025} documenta en la cadena de valor del valle Torobamba c\'omo la distribuci\'on desigual de ingresos deja a pequeños productores con menor capacidad de absorber ca\'idas de precio. La Figura~\ref{fig:formacion_precios} ilustra esquem\'aticamente la formaci\'on de precios y reparto de ingresos a lo largo de la cadena.

\incluirfigura{figuras/figura7\_2.jpg}{0.88\linewidth}{Formaci\'on de precios y distribuci\'on de ingresos en la cadena de palta}{fig:formacion_precios}

\incluirfigura{figuras/figura7\_3.jpg}{0.85\linewidth}{Ingresos totales por venta de palta Hass, a\~no 2024}{fig:ingresos_2024}

Cuando el precio FOB cae mientras los costos de insumos, agua, energ\'ia y mano de obra aumentan, el productor enfrenta m\'argenes negativos o break-even precarios. \citet{Velasquez2025} propone m\'etodos de cuantificaci\'on de riesgo de mercado aplicados al aguacate Hass, subrayando que la volatilidad no es un evento aislado sino una caracter\'istica estructural del rubro. Exportadores con contratos y diversificaci\'on pueden mitigar parcialmente el riesgo; pequeños productores desvinculados de esquemas contractuales permanecen m\'as expuestos.

Las Figuras~\ref{fig:precios_chile} y~\ref{fig:precios_usa} muestran referencias de precios mayoristas en Chile y Estados Unidos, mercados relevantes para benchmarking. Si bien no equivalen al FOB recibido por el productor peruano, reflejan la sensibilidad de precios spot a cambios de oferta y demanda en ventanas cortas.

\incluirfigura{figuras/figura2\_4.jpg}{0.88\linewidth}{Precios promedio por kilo de palta Hass en mercado mayorista Lo Valledor, Santiago de Chile, 2021}{fig:precios_chile}

\incluirfigura{figuras/figura2\_5.jpg}{0.92\linewidth}{Aguacate: precios medios al por mayor en EE.UU. (ene 2023 -- nov 2025)}{fig:precios_usa}

\subsection{Riesgo asociado a la toma de decisiones}

La decisi\'on de siembra de palta Hass implica un horizonte de retorno de mediano y largo plazo: desde la plantaci\'on hasta la producci\'on comercial pueden transcurrir cuatro a seis a\~nos, periodo durante el cual los precios internacionales pueden variar ampliamente \citep{VitteryFlores2023}. Un productor que expande hect\'areas bajo precios altos puede enfrentar m\'argenes reducidos al entrar en producci\'on, especialmente si otros productores tomaron decisiones similares basadas en la misma se\~nal observada.

La dependencia de experiencia emp\'irica y promedios hist\'oricos resulta insuficiente cuando la serie presenta quiebres estructurales ---como los observados tras la pandemia o ante la entrada acelerada de nuevos or\'igenes exportadores--- \citep{ReisFilho2022}. Sin herramientas predictivas, el productor sobrepondera el \'ultimo precio observado o relatos locales, incurriendo en sesgos de disponibilidad y subestimaci\'on del riesgo de ca\'ida.

El riesgo no es solo financiero: incluye sobreinversi\'on en infraestructura de riego, endeudamiento y presi\'on sobre recursos h\'idricos en valles ya saturados. Por ello, anticipar escenarios de precio FOB a horizontes de 4, 6 y 8 semanas ---\'util para decisiones operativas de campa\~na--- y combinarlos con proyecciones de rentabilidad constituye un aporte directo a la gesti\'on del riesgo en La Libertad y regiones palteras afines.
